\documentclass[
  % Babel language, also used to load translations
  german,
  % Use A4 paper size, you can change this to eg. letterpaper if you need
  % the letter format. The normal methods to modify the paper size should
  % be picked up by SDAPS automatically.
  % a4paper, % setting this might break the example scan unfortunately
  % letterpaper
  %
  % If you need it, you can add a custom barcode at the center
  globalid=Sommersemester-16,
  %
  % And the following adds a per sheet barcode at the bottom left
  print_questionnaire_id,
  %
  % You can choose between twoside and oneside. twoside is the default, and
  % requires the document to be printed and scanned in duplex mode.
  twoside,
  %
  % QR Codes instead Barcodes
  style=qr,
  % The following options make sense so that we can get a better feel for the
  % final look.
  pagemark,
  stamp]{sdaps}
\usepackage[utf8]{inputenc}
\usepackage{multicol}

%############ FRAGEBOGEN I ZRS ############%

%%% SCHRIFTARTEN %%%
% Überschriften auf Helvetica setzen.
\usepackage[scaled]{helvet}
% KOMAScript Schriftarten setzen.
\setkomafont{choicefont}{\sffamily}
\setkomafont{marklinequestionfont}{\sffamily}
% Schriftart global auf Sans Serif setzen.
\renewcommand{\familydefault}{\sfdefault}

%%% FARBEN %%%
% Farbe der Fragebogenabschnitte ändern
\definecolor{sectionbgcolor}{rgb}{0,0,0}
\definecolor{sectionfgcolor}{rgb}{255,255,255}

%%% LAYOUT ANPASSUNG %%%
% Trennstrich zwischen den einzelnen Fragen
\renewcommand{\separate}[1][0.12]{\vspace{#1cm}\noindent\makebox[\linewidth]{\rule{\textwidth}{0.5pt}}\vspace{#1cm}}


\author{Fabian Mundt}
\title{Zeit-Raum Studium --- Fragebogen Sommersemester 16}

\begin{document}
  % Everything you do should be done inside the questionnaire environment.

  % If you don't like the default text at the beginning of each questionnaire
  % you can remove it with the optional [noinfo] parameter for the environment
  \begin{questionnaire}[noinfo]
    % There is a predefined "info" style to hilight some text.


    %%%%%%%%%%%%%%%%%%%%%%%%%%%%%%%%%%%%%%%%%%%%%%%%%%%

    % PART 0 %
    \begin{center}
		\textbf{\textsc{\Huge Zeit-Raum Studium}} \\
		\vspace{.5cm}
		\textbf{\textsc{\LARGE Fragebogen Sommersemester 16}}
	\end{center}
	\vspace{.15cm}

	% Autoreninformationen löschen
	\thispagestyle{empty}
	% Seitenbegrenzungen setzen (für die maschinelle Verarbeitung besser)
	\sdapspagemark

    % INTRO (ALLGEMEINE HINWEISE)
    % !TEX root = Fragebogen_1.tex

% Metadaten für den Auswertungsbericht
\addinfo{Datum}{22.07.2015}
\addinfo{Befragungssituation}{Zweiter Pretest (Intensiver Pretest mit fünf Studierenden)}

% Introtext
\begin{info}
\vspace{.5cm}

Sehr geehrte Studentinnen und Studenten,

\vspace{.5em}

der folgende Fragebogen bezieht sich auf Ihren »Zeit-Raum Studium« an der Pädagogischen Hochschule Karlsruhe. Die Erhebung und Auswertung findet im Rahmen meiner Doktorarbeit an der Pädagogischen Hochschule Karlsruhe in der Allgemeinen und Historischen Erziehungswissenschaft statt. Geplant ist, \textbf{Sie in den ersten drei Semestern Ihres Bachelorstudiums forschend zu begleiten}. Dabei interessieren mich insbesondere Ihr biografischer Hintergrund, Ihre alltäglichen Abläufe und Ihre Studienerfahrungen. Ein tieferes Verständnis Ihres Studienalltags soll letztlich Ihnen zugutekommen, indem es ganz konkret auf die Verbesserung der Studienbedingungen vor Ort abzielt.

\vspace{.5em}

Der Fragebogen gliedert sich in \textbf{drei Teile}. Zunächst wird nach Ihrem \textbf{biografischen Hintergrund} gefragt, danach stehen Fragen zu Ihren Erfahrungen in der \textbf{Übergangszeit zwischen Schule und Hochschule} im Mittelpunkt. Den Abschluss bilden Fragen zu Ihren \textbf{alltäglichen Studienaktivitäten}.

\vspace{.5em}

Ich bitte Sie, sich Zeit zu nehmen und den Fragebogen ehrlich und vollständig zu bearbeiten. \textbf{Die Auswertung erfolgt anonym.} Da die Studie längsschnittlich (über drei Semester) angelegt ist und zudem einige Interviews durchgeführt werden sollen, wird nach Ihrer E-Mail-Adresse gefragt. Diese personenbezogene Angabe wird getrennt von Ihren übrigen Antworten aufbewahrt und  lediglich zur Kontaktaufnahme für eine Interviewanfrage herangezogen.

\vspace{.5em}

Vielen Dank im Voraus für Ihre Bereitschaft und Unterstützung!

\vspace{-1cm}

\includegraphics[width=0.25\textwidth]{unterschrift.png}

\vspace{-1.5cm}
\end{info}

\vspace{.75cm}

% Bearbeitungshinweise
\begin{info}
\vspace{.5em}

\textbf{\Large Bearbeitungshinweise}

Der Fragebogen wird maschinell erfasst. Um dabei optimale Ergebnisse zu erzielen, verwenden Sie zum Ausfüllen bitte einen \textbf{Kugelschreiber} oder einen \textbf{nicht zu dicken Filzstift} (es muss ein Kreuz erkennbar sein).

\vspace{.5em}

Wichtig ist, dass Sie ein \textbf{kräftiges (möglichst dunkles)} und \textbf{deutlich erkennbares Kreuz} setzen. Bitte verwenden Sie daher keine blassen Stifte (\textbf{keine Bleistifte}).

\vspace{.5em}

Beispiel für ein gut erfassbares Kreuz: {\LARGE \checkedbox{}}

\vspace{.5em}

Sie können einmalig Ihre Markierung korrigieren, indem Sie das Feld komplett ausfüllen.

\vspace{.5em}

Beispiel für eine Korrektur (Kreuz wird nicht erfasst): {\LARGE \correctedbox{}}

\vspace{.5em}

Sollten Sie Ihre Angaben mehrfach korrigieren wollen, notieren Sie ihre Angaben bitte handschriftlich.

\vspace{.5em}

\textbf{Bitte beachten Sie, dass der Fragebogen doppelseitig bedruckt ist. Bitte achten Sie auch darauf, die QR-Codes und die rechtwinkligen Seitenbegrenzungen nicht zu beschädigen.}

\end{info}


    \clearpage
    %%%%%%%%%%%%%%%%%%%%%%%%%%%%%%%%%%%%%%%%%%%%%%%%%%%
	\pagestyle{scrheadings}

	% PART III %
    \begin{center}
		\textbf{\textsc{\Huge Studienalltag}}
	\end{center}
    \vspace{.15cm}

	% PERSÖNLICHE SITUATION
	% !TEX root = ../Fragebogen_1.tex

% Persönliche Situation
\section{Persönliche Situation}
\vspace{.25cm}

\begin{choicequestion}[1]{Wo wohnen Sie \underline{während der Vorlesungszeit?}}
	\choiceitem{bei meinen Eltern}
	\choiceitem{in einer WG}
	\choiceitem{in einem Wohnheim}
	\choiceitem{allein in einer eigenen Wohnung}
	\choiceitem{mit meinem Partner/meiner Partnerin}
	\choiceitem{woanders}
\end{choicequestion}

\separate

\begin{choicequestion}[6]{Wie viele Kilometer sind es ungefähr von Ihrem Wohnort \underline{während der Vorlesungszeit} bis zur Hochschule?}
	\choiceitemtext{1cm}{1}{Kilometer:}
\end{choicequestion}

\separate

\begin{choicequestion}[6]{Wie lange brauchen Sie ungefähr von Ihrem Wohnort \underline{während der Vorlesungszeit} bis zur Hochschule?}
	\choiceitemtext{1cm}{1}{Minuten:~~~}
\end{choicequestion}

\separate

\begin{choicequestion}[1]{Wo wohnen Sie aller Voraussicht nach \underline{während der »Semesterferien«?}}
	\choiceitem{bei meinen Eltern}
	\choiceitem{in einer WG}
	\choiceitem{in einem Wohnheim}
	\choiceitem{allein in einer eigenen Wohnung}
	\choiceitem{mit meinem Partner/meiner Partnerin}
	\choiceitem{woanders}
	%\choiceitemtext{1cm}{1}{anderer Ort:}
\end{choicequestion}

\separate

\begin{choicequestion}[6]{Wie viele Kilometer sind es ungefähr von Ihrem Wohnort \underline{während der »Semesterferien«} bis zur Hochschule?}
	\choiceitemtext{1cm}{1}{Kilometer:}
\end{choicequestion}

\separate

\begin{choicequestion}[6]{Wie lange brauchen Sie ungefähr von Ihrem Wohnort \underline{während der »Semesterferien«} bis zur Hochschule?}
	\choiceitemtext{1cm}{1}{Minuten:~~~}
\end{choicequestion}

\separate

\begin{choicegroup}{Wie finanzieren Sie zur Zeit Ihr Studium?}
	\groupaddchoice{nein, dadurch nicht}
	\groupaddchoice{ja, teilweise}
	\groupaddchoice{ja, hauptsächlich}

	\choiceline{durch Unterstützung der Eltern}
	\choiceline{durch Einkommen des/ der (Ehe-)Partners/-in} %/der (Ehe-)Partnerin}
	\choiceline{durch BAföG}
	\choiceline{durch Studienkredit/ Bildungskredit}
	\choiceline{durch Stipendien (Förderwerke, Stiftungen,…)}
%	\choiceline{durch die Arbeit als studentische Hilfskraft}
	\choiceline{durch eigene Arbeit während der Vorlesungszeit}
	\choiceline{durch eigene Arbeit während der Semesterferien}
	\choiceline{durch anderes}
\end{choicegroup}

\vspace{-.1cm}
\separate
\vspace{-.1cm}

\textbox*{1cm}{Falls Sie einer bezahlten Arbeit nachgehen, um welche handelt es sich?}

\separate
\vspace{-.1cm}

\singlemarkLabelsSixLine{Wie kommen Sie insgesamt mit dem Geld zurecht, das Ihnen zur Verfügung steht?}{sehr schlecht}{ziemlich schlecht}{weder gut noch schlecht}{ziemlich gut}{sehr gut}{kann ich nicht sagen}

\vspace{-.25cm}
\separate
\vspace{-.15cm}

\begin{choicequestion}[6]{Wie viele Ihrer Kommiliton/innen würden Sie als Freund/innen bezeichnen?}
	\choiceitemtext{1cm}{1}{Anzahl:}
\end{choicequestion}

\separate
\vspace{-.15cm}

% Skala überdenken!!!
\singlemarkLabelsSixLine{Welchen Stellenwert nimmt Ihr Studium momentan in Ihrem Leben ein?}{unwichtig}{eher unwichtig}{teils teils}{eher wichtig}{sehr wichtig}{kann ich nicht sagen}

\vspace{-.25cm}
\separate
\vspace{-.25cm}

% AUF EINZELFRAGE UMSTELLEN
% Die Frage zielt auf die allgemeine Wahrnehmung der besuchten Schulen!
\begin{markgroupOne}{Wie würden Sie die PH Karlsruhe alles in allem tendenziell beschreiben?}
	\marklineRankingOne{modern}{altmodisch}
	\marklineRankingOne{neu}{abgenutzt}
	\marklineRankingOne{schön}{hässlich}
	\marklineRankingOne{offen}{geschlossen}
	\marklineRankingOne{groß}{klein}
	\marklineRankingOne{einladend}{abweisend}
	\marklineRankingOne{sauber}{schmutzig}
	\marklineRankingOne{monoton}{vielfältig}
\end{markgroupOne}%


	% EINE GWÖHNLICHE STUDIENWOCHE
	% !TEX root = Fragebogen_1.tex

\pagebreak
% gewöhnliche Studienwoche
\vspace{-.25cm}
\section{Gewöhnliche Studienwoche}

\question{Wie viel Zeit verbringen Sie pro Wochentag mit den angegebenen Tätigkeiten?
\newline\small{Bitte \underline{verteilen Sie die 24 Stunden eines Tages}. Beachten Sie das von Ihren Mentor/innen eingeblendete Beispiel.}}
% MIT SDAPS AUTOMATISIEREN!!!!
% >{\centering\arraybackslash} Array Syntax um Spalten zu formatieren
\begin{tabular}{|m{6.1cm}|*{7}{>{\centering\bf\arraybackslash}p{1.25cm}|}}

\hline
							& Mo 	& Di	& Mi 	& Do 	& Fr 	& Sa 	& So \\
% \hline
% \vspace{.50cm}
% \textbf{Zeit an der Hochschule} %\newline\footnotesize{(Wie lange sind auf dem Hochschulgelände)}
% \vspace{.50cm} 				& 		&		& 		& 		& 		&		& 	 \\
\hline
\vspace{.50cm}
\textbf{Lehrveranstaltungen} %\newline\footnotesize{(Vorlesungen, Seminare etc.)}
\vspace{.50cm} 				& 		&		& 		& 		& 		&		& 	 \\
\hline
\vspace{.50cm}
\textbf{Zeit zwischen den Veranstaltungen}\nolinebreak %\newline\footnotesize{(»Wartezeiten«)}
\vspace{.50cm} 				& 		&		& 		& 		& 		&		& 	 \\
\hline
\vspace{.50cm}
\textbf{Selbststudium} % \newline\footnotesize{(Lernen für das Studium)}
\vspace{.50cm} 				& 		&		& 		& 		& 		&		& 	 \\
\hline
\vspace{.50cm}
\textbf{Fahrtzeit}% \newline\footnotesize{(Zeit zur PH und zurück zur Wohnung.)}
\vspace{.50cm} 				& 		&		& 		& 		& 		&		& 	 \\
\hline
\vspace{.50cm}
\textbf{Arbeitszeit}% \newline\footnotesize{(Zeit zum Geld verdienen.)}
\vspace{.50cm} 				& 		&		& 		& 		& 		&		& 	 \\
\hline
\vspace{.50cm}
\textbf{Freizeit} %\newline\footnotesize{(Zeit für Ihre Hobbies etc.)}
\vspace{.50cm} 				& 		&		& 		& 		& 		&		& 	 \\
\hline
\vspace{.50cm}
\textbf{Schlafen}
\vspace{.50cm} 				& 		&		& 		& 		& 		&		& 	 \\
\hline
% \vspace{.50cm}
% \textbf{Zeit im eigenen Wohnbereich} %\newline\footnotesize{~}
% \vspace{.50cm} 				& 		&		& 		& 		& 		&		& 	 \\
% \hline
\end{tabular}
\vspace{.15cm}

% \separate
%
% \begin{choicegroupWide}{Wie nutzen Sie die Wartezeiten zwischen einzelnen Lehrveranstaltungen in der Regel?}
% 	\groupaddchoice{sehr\linebreak oft}
% 	\groupaddchoice{oft}
% 	\groupaddchoice{manch-\linebreak mal}
%   \groupaddchoice{selten}
% 	\groupaddchoice{sehr selten}
% 	\groupaddchoice{nie}
%
% 	\choicelineWide{Veranstaltungen vor- und/oder nachbereiten}
% 	\choicelineWide{Selbststudium}
% 	\choicelineWide{gemeinsam mit Kommiliton/innen lernen}
% 	\choicelineWide{essen und trinken}
% 	\choicelineWide{mit Freunden treffen}
%   	\choicelineWide{shoppen}
% 	\choicelineWide{entspannen}
% 	\choicelineWide{Haushaltseinkäufe erledigen}
% 	\choicelineWide{im Internet surfen}
% 	\choicelineWide{in die Wohnung gehen}
%
% \end{choicegroupWide}


% \textbox*{3cm}{Wie nutzen Sie die Wartezeiten zwischen einzelnen Lehrveranstaltungen in der Regel?}

\separate

\begin{choicequestion}[6]{An wie vielen Lehrveranstaltungen nehmen Sie dieses Semester teil?}
	\choiceitemtext{1cm}{1}{Anzahl:}
\end{choicequestion}

\separate

\begin{choicequestion}[1]{Sind Sie alles in allem mit Ihrem Lehrveranstaltungs\underline{plan} zufrieden?}
	\choiceitem{ja}
	\choiceitem{teilweise}
	\choiceitem{nein}
\end{choicequestion}

\separate

\begin{choicequestion}[1]{Mit welchem Fortbewegungsmittel kommen Sie hauptsächlich zur Hochschule?}
	\choiceitem{Auto}
	\choiceitem{Motorrad/Roller}
	\choiceitem{Bus}
	\choiceitem{Bahn}
	\choiceitem{Fahrrad}
	\choiceitem{zu Fuß}
\end{choicequestion}


	% GEWÖHNLICHE STUDIENTÄTIGKEITEN INNERHALB EINER WOCHE
	% !TEX root = Fragebogen_1.tex

% alltägliche Studienaktivitäten
\vspace{-.25cm}
\section{Alltägliche Studienaktivitäten}
\vspace{-.15cm}

\begin{choicegroupWide}{Welchen Tätigkeiten gehen Sie in einer typischen Studienwoche \underline{an der PH Karlsruhe} nach?}
	\groupaddchoice{sehr\linebreak oft}
	\groupaddchoice{oft}
	\groupaddchoice{manch-\linebreak mal}
	\groupaddchoice{selten}
	\groupaddchoice{sehr selten}
	\groupaddchoice{nie}

	% PH %%%%%%%%%%%%%%%%%%%%%%%%%%%%%%%%%%%%%%%%%%%%%%%%%%%%%%%%%%%%%%%%%%%%

	% PH Bibliothek
	\choicelineWide{in der PH Bibliothek (HSB) recherchieren}
	\choicelineWide{in der PH Bibliothek (HSB) lernen}
	% Mensa
	\choicelineWide{in der Mensa essen/trinken}
	% Cafeteria
	\choicelineWide{in der Cafeteria essen/trinken}
	\choicelineWide{in der Cafeteria lernen}
	% Foyer
	\choicelineWide{im Foyer lernen}
	% Lernräume
	\choicelineWide{in den Lern- und Übungsräumen lernen}
	% Lerninseln
	\choicelineWide{die Lerninseln zum Lernen nutzen}
	% Computerräume
	\choicelineWide{die Computerräume zum Studieren nutzen}
	% Kunstwerkstätten
	\choicelineWide{die (Kunst)Werkstätten nutzen}
	% % ZIM
	\choicelineWide{Angebote des Medienzentrums nutzen}
	% % SSZ
	% \choicelineWide{Angebote des SSZ nutzen}
	% % SSZ
	% \choicelineWide{Angebote des SLZ nutzen}
	% % Lernzentrum
	\choicelineWide{Angebote des Lehr-Lern-Zentrums nutzen}
	% % Medienlabor
	% \choicelineWide{das Medienlabor (in Gebäude 1) nutzen}
	% Veranstaltungsräume
	\choicelineWide{in leerstehenden Veranstaltungsräumen lernen}
	% Campus lernen
	\choicelineWide{auf dem Campus lernen}
	% Campus entspannen
	\choicelineWide{auf dem Campus entspannen}
	% KOMMILITONEN %%%%%%%%%%%%%%%%%%%%%%%%%%%%%%%%%%%%%%%%%%%%%%%%%%%%%%%%%%
	% Fachschaft
	\choicelineWide{in einer Fachschaft engagieren}
	% StuVe
	\choicelineWide{in der StuVe engagieren}
	% Lehrende
	\choicelineWide{mit Lehrenden austauschen (nicht in Veranstalt.)}
	% Kommilitonen
	\choicelineWide{mit Kommiliton/innen entspannen}
	% Kommilitonen
	\choicelineWide{mit Kommiliton/innen zum Lernen treffen}
\end{choicegroupWide}

% \separate

% \textbox*{1.5cm}{Falls Sie nicht genannten Tätigkeiten innerhalb der PH Karlsruhe an nicht genannten Orten nachgehen, welche sind das?}

\vspace{-.27cm}
\separate
\vspace{-.27cm}

\begin{choicegroupWide}{Welchen Tätigkeiten gehen Sie in einer typischen Studienwoche \underline{außerhalb der PH} nach?}
	\groupaddchoice{sehr\linebreak oft}
	\groupaddchoice{oft}
	\groupaddchoice{manch-\linebreak mal}
	\groupaddchoice{selten}
	\groupaddchoice{sehr selten}
	\groupaddchoice{nie}

	% ANDERE EINRICHTUNGEN %%%%%%%%%%%%%%%%%%%%%%%%%%%%%%%%%%%%%%%%%%%%%%%%%%

	% Bibliotheken
	\choicelineWide{in anderen Bibliotheken (z.B. BLB) lernen}
	% Einrichtungen anderer Unis
	\choicelineWide{andere öffentliche Einrichtungen (z.B. KIT) zum Lernen nutzen}

	% STADT KARLSRUHE %%%%%%%%%%%%%%%%%%%%%%%%%%%%%%%%%%%%%%%%%%%%%%%%%%%%%%%

	% Stadt lernen
	\choicelineWide{in der Stadt lernen}
	% Stadt essen
	\choicelineWide{in der Stadt essen}
	% Stadt entspannen
	\choicelineWide{in der Stadt entspannen}
	% Cafés
	\choicelineWide{Cafés zum Lernen nutzen}
	% Kneipen, Pubs
	\choicelineWide{Kneipen, Pubs etc. zum Lernen nutzen}
	% Parks
	\choicelineWide{Parkanlagen zum Lernen nutzen}
	% Bahn
	\choicelineWide{In der Bahn lernen}

	% WOHNUNG %%%%%%%%%%%%%%%%%%%%%%%%%%%%%%%%%%%%%%%%%%%%%%%%%%%%%%%%%%%%%%%

	% Eigene Wohnung
	\choicelineWide{am eigenen Wohnort lernen}
	% Fremde Wohnung
	\choicelineWide{an einem fremden Wohnort (z.B. Freunde) lernen}

	% KOMMILITONEN %%%%%%%%%%%%%%%%%%%%%%%%%%%%%%%%%%%%%%%%%%%%%%%%%%%%%%%%%%

	% Kommilitonen
	\choicelineWide{mit Kommiliton/innen zum Entspannen treffen}
	% Kommilitonen
	\choicelineWide{mit Kommiliton/innen zum Lernen treffen}
\end{choicegroupWide}

% \separate
%
% \textbox*{1.5cm}{Falls Sie nicht genannten Tätigkeiten außerhalb der PH Karlsruhe an nicht genannten Orten nachgehen, welche sind das?}

\vspace{-.36cm}
\separate
\vspace{-.25cm}

\textbox*{1cm}{Wo ist Ihr bevorzugter Lernort?}

\separate

\textbox*{2cm}{Was sind Ihre Gründe für die Wahl dieses Lernorts?}

\separate

\begin{choicequestion}[1]{Nehmen Sie an einer Lerngruppe teil?}
	\choiceitem{ja}
	\choiceitem{gelegentlich}
	\choiceitem{nein}
\end{choicequestion}

\separate

\textbox*{1cm}{Falls Sie an einer Lerngruppe teilnehmen, wo ist Ihr bevorzugter Lernort?}

\separate

\textbox*{2cm}{Falls Sie an einer Lerngruppe teilnehmen, was sind Ihre Gründe für die Wahl dieses Lernorts?}

\separate

\begin{choicegroupWide}{Wie nutzen Sie das Internet in einer gewöhnlichen Studienwoche \underline{für Ihr Studium?}}
	\groupaddchoice{sehr\linebreak oft}
	\groupaddchoice{oft}
	\groupaddchoice{manch-\linebreak mal}
	\groupaddchoice{selten}
	\groupaddchoice{sehr selten}
	\groupaddchoice{nie}

	% Services der Hochschuel
	\choicelineWide{Website der Hochschule}

	\choicelineWide{Web Services der PH (z.B. Stud.IP, LSF)}

	\choicelineWide{Andere Web Services (z.B. YouTube, MOOCs)}

	\choicelineWide{E-Mails an Lehrpersonen schreiben}

	\choicelineWide{E-Mails an Kommiliton/innen schreiben}

	\choicelineWide{für das Studium recherchieren}

	\choicelineWide{Soziale Netzwerke (z.B. spezielle Gruppen)}

	\choicelineWide{Nachrichtendienste (z.B. WhatsApp)}

\end{choicegroupWide}

% \separate
%
% \textbox*{3cm}{Falls Sie nicht genannten Studientätigkeiten im Internet nachgehen, welche sind das?}

\separate

\singlemarkLabelsSixLine{Zu welcher Tageszeit sind Sie am aufmerksamsten für Studienzwecke?}{morgens}{mittags}{nachmittags}{abends}{nachts}{kann ich\linebreak nicht sagen}

\vspace{-.55cm}
\separate
\vspace{-.5cm}

\singlemarkLabelsSixLine{Wo ist Ihr bevorzugter Sitzplatz in einer gewöhnlichen Veranstaltung?}{in der ersten Reihe}{im vorderen Drittel}{in der Mitte}{im hinteren Drittel}{in der letzten Reihe}{kann ich\linebreak nicht sagen}

\vspace{-.55cm}
\separate
\vspace{-.5cm}

\singlemarkLabelsSixLine{Wählen Sie, falls möglich, in jeder Veranstaltung den selben Sitzplatz?}{nie}{\hspace{10pt}selten}{\hspace{5pt}ab und zu}{\hspace{10pt}oft}{\hspace{10pt}immer}{kann ich \linebreak nicht sagen}

\vspace{-.55cm}
\separate
\vspace{-.5cm}

% Anzahl benennen?
\singlemarkLabelsSixLine{Wie oft nutzen Sie Ihr Smartphone (für jegliche Zwecke) während einer gewöhnlichen Veranstaltung?}{nie}{\hspace{10pt}selten\newline \footnotesize{(1 bis 5 mal)}}{\hspace{5pt}ab und zu\newline \footnotesize{(5 bis 15 mal)}}{\hspace{10pt}oft\newline \footnotesize{(15 bis 30 mal)}}{\hspace{10pt}immer\newline \footnotesize{(über 30 mal)}}{kann ich \linebreak nicht sagen}

% \vspace{-.5cm}
% \separate
% \vspace{-.5cm}
%
% DIESE FRAGE IN DEN LEITFADEN AUFNEHMEN
%
% \singlemarkLabelsSixLine{Wie empfinden Sie die in der Regel 90-minütige Dauer einer Veranstaltung?}{sehr schlecht}{ziemlich schlecht}{weder gut noch schlecht}{ziemlich gut}{sehr gut}{kann ich\linebreak nicht sagen}


	% GEWÖHNLICHE FREIZEITAKTIVITÄTEN INNERHALB EINER WOCHE
	% !TEX root = Fragebogen_2.tex

\vspace{-.6cm}
\section{Alltägliche Freizeitaktivitäten}
\vspace{-.2cm}

\begin{choicegroupWide}{Welchen Freizeittätigkeiten gehen Sie in einer gewöhnlichen Woche nach?}
	\groupaddchoice{sehr\linebreak oft}
	\groupaddchoice{oft}
	\groupaddchoice{manch-\linebreak mal}
  	\groupaddchoice{selten}
	\groupaddchoice{sehr selten}
	\groupaddchoice{nie}

	\choicelineWide{Fernsehen (auch DVD/Blu-Ray/Streaming)}
	\choicelineWide{Musik hören}
	% \choicelineWide{Videos/DVDs anschauen}
	\choicelineWide{ins Kino gehen}
	\choicelineWide{sich mit Freunden treffen}
	\choicelineWide{chatten (z.B. WhatsApp)}
	\choicelineWide{Soziale Netzwerke (z.B. Facebook) nutzen}
	\choicelineWide{im Internet surfen}
	\choicelineWide{nichts tun, ›rumhängen‹}
	\choicelineWide{Bücher lesen}
	% \choicelineWide{in die Kneipe gehen}
	% \choicelineWide{Zeitschriften oder Magazine lesen}
	\choicelineWide{Zeitungen lesen}
	\choicelineWide{in die Disco, zu Partys oder Konzerten gehen}
	\choicelineWide{Computerspiele spielen}
	\choicelineWide{Sport treiben (Freizeit, Verein…)}
	% \choicelineWide{Sport in der Freizeit, wie Rad fahren, Skaten, Kicken usw.}
	% \choicelineWide{Training/ aktiv Sport treiben (Fitnessclub, Sportverein ...)}
	% \choicelineWide{sich mit Leuten (keine Kommilitonen) treffen}
	% \choicelineWide{sich mit Kommilitonen treffen}
	% \choicelineWide{sich in einem Projekt/ einer Initiative/ einem Verein engagieren}
	\choicelineWide{etwas mit der Familie unternehmen}
	\choicelineWide{shoppen, sich tolle Sachen kaufen}
	\choicelineWide{etwas Kreatives, Künstlerisches machen}
	% \choicelineWide{Oper, klassische Konzerte, Theater besuchen}
	\choicelineWide{Kunstausstellungen, Museen besuchen}
	\choicelineWide{religiöse Veranstaltungen (z.B. Kirche) besuchen}
	\choicelineWide{ehrenamtlich engagieren}
	\choicelineWide{politisch engagieren}
	% \choicelineWide{Tagesausflüge/ kleinere Reisen unternehmen}
	% \choicelineWide{in die Natur gehen}
	% \choicelineWide{in die Stadt gehen}
	% \choicelineWide{in die Natur gehen}
	% \choicelineWide{Cafés besuchen}
	% \choicelineWide{in die Natur gehen}
	% \choicelineWide{in Restaurants essen gehen}
\end{choicegroupWide}
\vspace{-5cm}


    %%%%%%%%%%%%%%%%%%%%%%%%%%%%%%%%%%%%%%%%%%%%%%%%%%%

    \newpage
    % PART IV %
    \begin{center}
		\textbf{\textsc{\Huge Schluss}}
	\end{center}
	\vspace{.15cm}

	% EXTRO (IDENTIFIKATION UND ANONYMISIERUNGSHINWEISE)
	% !TEX root = Fragebogen_1.tex

% Extro
\vspace{-.25cm}
\section{Weitere Teilnahme an der Studie}
\vspace{.25cm}

\begin{choicequestion}[4]{Die Studie versucht Sie im Rahmen der ersten drei Studiensemester zu begleiten. Um künftige Befragungen zuordnen zu können, wird ein anonymisierter Zuordnungsschlüssel benötigt.}
	\choiceitemtext{1cm}{1}{\makebox[11cm]{Erste \underline{zwei Buchstaben} des Vornamens Ihrer \underline{Mutter} (z.B. \textit{JU} für \textit{\underline{Ju}lia}):\hfill}}\\
	\choiceitemtext{1cm}{1}{\makebox[11cm]{Erste \underline{zwei Buchstaben} des Vornamens Ihres \underline{Vaters} (z.B. \textit{PA} für \textit{\underline{Pa}ul}):\hfill}}\\
	\choiceitemtext{1cm}{1}{\hspace*{-0.2cm}\makebox[11.2cm]{Der \underline{Tag} Ihres Geburtsdatums (z.B. \textit{17} für den \textit{\underline{17}.09.1997}):\hfill}}
\end{choicequestion}

\separate

\begin{choicequestion}[1]{Wären Sie bereit, sich interviewen zu lassen?}
	\choiceitem{ja}
	\choiceitem{nein}
	\choiceitemtext{1cm}{1}{Falls ja, Ihre E-Mail-Adresse:}
\end{choicequestion}

\vspace{.25cm}
\section{Rückmeldung}
\vspace{.25cm}


\textbox*{5.5cm}{Wenn Sie noch Ergänzungen, Anmerkungen oder kritische Hinweise haben, bitte ich Sie, mir diese hier mitzuteilen.}

\separate

\textbf{Vielen Dank, dass Sie mich durch das Ausfüllen des Fragebogens bei meiner Doktorarbeit unterstützt haben!}

\vspace{-1.5cm}
\flushright
\includegraphics[width=0.25\textwidth]{unterschrift.png}
\vspace{-1.5cm}
\flushright
\href{mailto:fabian.mundt@ph-karlsruhe.de}{\nolinkurl{fabian.mundt@ph-karlsruhe.de}}


    %%%%%%%%%%%%%%%%%%%%%%%%%%%%%%%%%%%%%%%%%%%%%%%%%%%

  \end{questionnaire}
\end{document}
